% frontpage.tex (METADATA ONLY)
% Use with:
%   % frontpage.tex (METADATA ONLY)
% Use with:
%   % frontpage.tex (METADATA ONLY)
% Use with:
%   % frontpage.tex (METADATA ONLY)
% Use with:
%   \input{frontpage}
%   \beforepreface
%   \afterpreface

% -------------------------
% REQUIRED METADATA
% -------------------------
\title{PARTIAL DISCHARGE CLASSIFICATION USING MACHINE LEARNING}
\author{MALCOLM SABINUS}

\degree{BACHELOR OF SCIENCE WITH HONORS INDUSTRIAL PHYSICS}
\dept{INDUSTRIAL PHYSICS}
\mainsupervisor{Dr. Rodney Petrus Balandong}
\matricno{BS22110396}

% Submission date format (MONTH YEAR)
\submitdate{%
\ifcase\the\month\or
JANUARY\or FEBRUARY\or MARCH\or APRIL\or MAY\or JUNE\or
JULY\or AUGUST\or SEPTEMBER\or OCTOBER\or NOVEMBER\or DECEMBER\fi
\space \number\the\year}

\copyrightyear{\number\the\year}

% NOTE: Avoid \hline here. It's table-only and can break formatting.
% If you need separators, use plain \\ only.
\authoraddress{%
1234 Elm Street, Apt 567\\
Springfield, 98765\\
Sabah, Malaysia%
}

% -------------------------
% DECLARATION TEXT (inside the box on the Declaration page)
% -------------------------
\declarationLLM{%
I acknowledge that during the preparation of this thesis, I have utilized Large
Language Models (LLMs) as a supplementary tool to assist in various aspects of
my research and writing process. Specifically, I have used LLMs for paraphrasing
and refining my writing to improve clarity, coherence, and readability while
ensuring the original meaning remains intact. Additionally, I have leveraged LLMs
to generate general ideas and structure my thoughts more effectively, using them
as a brainstorming aid rather than as a replacement for original critical thinking.
LLMs have also been helpful in summarizing complex topics, identifying key points
from lengthy articles, and suggesting alternative ways to present information.
Furthermore, I have used LLMs to check grammar and syntax, ensuring that my
writing adheres to academic standards. In some instances, LLMs provided
guidance on citation formats and referencing best practices, although all sources
have been manually verified to maintain academic integrity. Importantly, I have
critically evaluated all AI-generated content and ensured that it aligns with my
own understanding and research findings.%
}

% -------------------------
% ACKNOWLEDGEMENT (optional: keep short, 1 page max per guide)
% If you don't want it printed, you can leave it empty.
% -------------------------
\penghargaan{\input{penghargaan.tex}
}

% -------------------------
% ABSTRACT (English) - 150 to 250 words, single paragraph, single spacing
% fin.sty will print it; you only store the text here.
% -------------------------
\abstract{Partial discharge (PD) is a critical indicator of insulation degradation in high-voltage electrical equipment, and accurate PD classification is essential for condition-based maintenance and early fault prevention. Conventional PD diagnostic approaches rely heavily on expert interpretation of phase-resolved discharge patterns, which are subjective, time-consuming, and vulnerable to noise and complex operating conditions. To address these limitations, this study investigates the use of traditional machine learning techniques for automated partial discharge classification using a structured, multi-domain feature-based framework. A publicly available partial discharge dataset was utilized, and features were extracted from time-domain, statistical, frequency-domain, and time--frequency representations to capture the physical and spectral characteristics of PD signals. Two independent feature selection strategies---Featurewiz and MLJAR AutoML---were employed to identify the most informative features while reducing dimensionality. Several traditional machine learning classifiers, including Support Vector Machine, Random Forest, Decision Tree, k-Nearest Neighbors, and Artificial Neural Network models, were trained and evaluated using stratified k-fold cross-validation. Model performance was assessed using accuracy, precision, recall, and F1-score metrics. The results show that Random Forest and Artificial Neural Network classifiers consistently achieved the strongest performance across both feature selection strategies, with classification accuracies exceeding 90\% and stable F1-scores. Despite differences in feature selection mechanisms, both approaches identified a common core of discriminative features, particularly those related to phase-resolved discharge behavior and energy distribution across frequency bands. Overall, this study confirms the practical viability of traditional machine learning approaches for robust, interpretable, and computationally efficient partial discharge classification in real-world power system monitoring applications.
}

% -------------------------
% ABSTRACT (Bahasa Malaysia) - required after English abstract
% -------------------------
\abstractMalay{Pelepasan separa (Partial Discharge, PD) merupakan petunjuk awal yang kritikal terhadap degradasi penebatan dalam peralatan elektrik voltan tinggi, dan pengelasan PD yang tepat adalah penting bagi penyelenggaraan berasaskan keadaan serta pencegahan kegagalan awal. Kaedah diagnostik PD konvensional bergantung secara meluas kepada tafsiran pakar terhadap corak pelepasan terlerai fasa, yang bersifat subjektif, memakan masa, dan terdedah kepada gangguan hingar serta keadaan operasi yang kompleks. Bagi mengatasi kekangan ini, kajian ini menilai penggunaan teknik pembelajaran mesin tradisional untuk pengelasan pelepasan separa secara automatik menggunakan rangka kerja ciri berbilang domain yang tersusun. Set data pelepasan separa yang tersedia secara umum telah digunakan, dan ciri-ciri diekstrak daripada domain masa, statistik, domain frekuensi, dan masa--frekuensi bagi menangkap ciri fizikal dan spektrum isyarat PD. Dua strategi pemilihan ciri yang bebas, iaitu Featurewiz dan MLJAR AutoML, telah diaplikasikan bagi mengenal pasti ciri paling diskriminatif serta mengurangkan dimensi ciri. Beberapa pengelas pembelajaran mesin tradisional, termasuk Support Vector Machine, Random Forest, Decision Tree, k-Nearest Neighbors, dan Artificial Neural Network, telah dilatih dan dinilai menggunakan pengesahan silang k-lipatan berstrata. Prestasi model dinilai berdasarkan metrik ketepatan, kejituan, kebolehan panggilan semula, dan skor F1. Keputusan menunjukkan bahawa pengelas Random Forest dan Artificial Neural Network secara konsisten mencapai prestasi terbaik merentas kedua-dua strategi pemilihan ciri, dengan ketepatan pengelasan melebihi 90\% dan skor F1 yang stabil. Walaupun terdapat perbezaan dalam mekanisme pemilihan ciri, kedua-dua pendekatan mengenal pasti satu set teras ciri yang signifikan, khususnya yang berkaitan dengan tingkah laku pelepasan terlerai fasa dan taburan tenaga merentasi jalur frekuensi. Secara keseluruhan, kajian ini mengesahkan kebolehlaksanaan praktikal kaedah pembelajaran mesin tradisional untuk pengelasan pelepasan separa yang teguh, boleh ditafsir, dan cekap dari segi pengiraan dalam aplikasi pemantauan sistem kuasa sebenar.
}

%   \beforepreface
%   \afterpreface

% -------------------------
% REQUIRED METADATA
% -------------------------
\title{PARTIAL DISCHARGE CLASSIFICATION USING MACHINE LEARNING}
\author{MALCOLM SABINUS}

\degree{BACHELOR OF SCIENCE WITH HONORS INDUSTRIAL PHYSICS}
\dept{INDUSTRIAL PHYSICS}
\mainsupervisor{Dr. Rodney Petrus Balandong}
\matricno{BS22110396}

% Submission date format (MONTH YEAR)
\submitdate{%
\ifcase\the\month\or
JANUARY\or FEBRUARY\or MARCH\or APRIL\or MAY\or JUNE\or
JULY\or AUGUST\or SEPTEMBER\or OCTOBER\or NOVEMBER\or DECEMBER\fi
\space \number\the\year}

\copyrightyear{\number\the\year}

% NOTE: Avoid \hline here. It's table-only and can break formatting.
% If you need separators, use plain \\ only.
\authoraddress{%
1234 Elm Street, Apt 567\\
Springfield, 98765\\
Sabah, Malaysia%
}

% -------------------------
% DECLARATION TEXT (inside the box on the Declaration page)
% -------------------------
\declarationLLM{%
I acknowledge that during the preparation of this thesis, I have utilized Large
Language Models (LLMs) as a supplementary tool to assist in various aspects of
my research and writing process. Specifically, I have used LLMs for paraphrasing
and refining my writing to improve clarity, coherence, and readability while
ensuring the original meaning remains intact. Additionally, I have leveraged LLMs
to generate general ideas and structure my thoughts more effectively, using them
as a brainstorming aid rather than as a replacement for original critical thinking.
LLMs have also been helpful in summarizing complex topics, identifying key points
from lengthy articles, and suggesting alternative ways to present information.
Furthermore, I have used LLMs to check grammar and syntax, ensuring that my
writing adheres to academic standards. In some instances, LLMs provided
guidance on citation formats and referencing best practices, although all sources
have been manually verified to maintain academic integrity. Importantly, I have
critically evaluated all AI-generated content and ensured that it aligns with my
own understanding and research findings.%
}

% -------------------------
% ACKNOWLEDGEMENT (optional: keep short, 1 page max per guide)
% If you don't want it printed, you can leave it empty.
% -------------------------
\penghargaan{First and foremost, I would like to express my sincere gratitude to my supervisor, Dr. Rodney Petrus Balandong, for his invaluable guidance, continuous support, and constructive feedback throughout the course of this research. His expertise,  patience, and dedication to academic excellence greatly contributed to the successful completion of this thesis.
I would also like to acknowledge the Faculty of Science and Technology, Universiti Malaysia Sabah, for providing the academic environment, facilities, and resources necessary to conduct this study. The support and assistance from the faculty members and administrative staff are sincerely appreciated.
My deepest appreciation is extended to my family for their unwavering encouragement, understanding, and moral support throughout my academic journey. Their constant motivation and belief in me have been a source of strength during challenging times.
I would also like to thank my friends and peers who provided helpful discussions, technical assistance, and encouragement throughout this research. Their support contributed positively to the completion of this work.
}

% -------------------------
% ABSTRACT (English) - 150 to 250 words, single paragraph, single spacing
% fin.sty will print it; you only store the text here.
% -------------------------
\abstract{Partial discharge (PD) is a critical indicator of insulation degradation in high-voltage electrical equipment, and accurate PD classification is essential for condition-based maintenance and early fault prevention. Conventional PD diagnostic approaches rely heavily on expert interpretation of phase-resolved discharge patterns, which are subjective, time-consuming, and vulnerable to noise and complex operating conditions. To address these limitations, this study investigates the use of traditional machine learning techniques for automated partial discharge classification using a structured, multi-domain feature-based framework. A publicly available partial discharge dataset was utilized, and features were extracted from time-domain, statistical, frequency-domain, and time--frequency representations to capture the physical and spectral characteristics of PD signals. Two independent feature selection strategies---Featurewiz and MLJAR AutoML---were employed to identify the most informative features while reducing dimensionality. Several traditional machine learning classifiers, including Support Vector Machine, Random Forest, Decision Tree, k-Nearest Neighbors, and Artificial Neural Network models, were trained and evaluated using stratified k-fold cross-validation. Model performance was assessed using accuracy, precision, recall, and F1-score metrics. The results show that Random Forest and Artificial Neural Network classifiers consistently achieved the strongest performance across both feature selection strategies, with classification accuracies exceeding 90\% and stable F1-scores. Despite differences in feature selection mechanisms, both approaches identified a common core of discriminative features, particularly those related to phase-resolved discharge behavior and energy distribution across frequency bands. Overall, this study confirms the practical viability of traditional machine learning approaches for robust, interpretable, and computationally efficient partial discharge classification in real-world power system monitoring applications.
}

% -------------------------
% ABSTRACT (Bahasa Malaysia) - required after English abstract
% -------------------------
\abstractMalay{Pelepasan separa (Partial Discharge, PD) merupakan petunjuk awal yang kritikal terhadap degradasi penebatan dalam peralatan elektrik voltan tinggi, dan pengelasan PD yang tepat adalah penting bagi penyelenggaraan berasaskan keadaan serta pencegahan kegagalan awal. Kaedah diagnostik PD konvensional bergantung secara meluas kepada tafsiran pakar terhadap corak pelepasan terlerai fasa, yang bersifat subjektif, memakan masa, dan terdedah kepada gangguan hingar serta keadaan operasi yang kompleks. Bagi mengatasi kekangan ini, kajian ini menilai penggunaan teknik pembelajaran mesin tradisional untuk pengelasan pelepasan separa secara automatik menggunakan rangka kerja ciri berbilang domain yang tersusun. Set data pelepasan separa yang tersedia secara umum telah digunakan, dan ciri-ciri diekstrak daripada domain masa, statistik, domain frekuensi, dan masa--frekuensi bagi menangkap ciri fizikal dan spektrum isyarat PD. Dua strategi pemilihan ciri yang bebas, iaitu Featurewiz dan MLJAR AutoML, telah diaplikasikan bagi mengenal pasti ciri paling diskriminatif serta mengurangkan dimensi ciri. Beberapa pengelas pembelajaran mesin tradisional, termasuk Support Vector Machine, Random Forest, Decision Tree, k-Nearest Neighbors, dan Artificial Neural Network, telah dilatih dan dinilai menggunakan pengesahan silang k-lipatan berstrata. Prestasi model dinilai berdasarkan metrik ketepatan, kejituan, kebolehan panggilan semula, dan skor F1. Keputusan menunjukkan bahawa pengelas Random Forest dan Artificial Neural Network secara konsisten mencapai prestasi terbaik merentas kedua-dua strategi pemilihan ciri, dengan ketepatan pengelasan melebihi 90\% dan skor F1 yang stabil. Walaupun terdapat perbezaan dalam mekanisme pemilihan ciri, kedua-dua pendekatan mengenal pasti satu set teras ciri yang signifikan, khususnya yang berkaitan dengan tingkah laku pelepasan terlerai fasa dan taburan tenaga merentasi jalur frekuensi. Secara keseluruhan, kajian ini mengesahkan kebolehlaksanaan praktikal kaedah pembelajaran mesin tradisional untuk pengelasan pelepasan separa yang teguh, boleh ditafsir, dan cekap dari segi pengiraan dalam aplikasi pemantauan sistem kuasa sebenar.
}

%   \beforepreface
%   \afterpreface

% -------------------------
% REQUIRED METADATA
% -------------------------
\title{PARTIAL DISCHARGE CLASSIFICATION USING MACHINE LEARNING}
\author{MALCOLM SABINUS}

\degree{BACHELOR OF SCIENCE WITH HONORS INDUSTRIAL PHYSICS}
\dept{INDUSTRIAL PHYSICS}
\mainsupervisor{Dr. Rodney Petrus Balandong}
\matricno{BS22110396}

% Submission date format (MONTH YEAR)
\submitdate{%
\ifcase\the\month\or
JANUARY\or FEBRUARY\or MARCH\or APRIL\or MAY\or JUNE\or
JULY\or AUGUST\or SEPTEMBER\or OCTOBER\or NOVEMBER\or DECEMBER\fi
\space \number\the\year}

\copyrightyear{\number\the\year}

% NOTE: Avoid \hline here. It's table-only and can break formatting.
% If you need separators, use plain \\ only.
\authoraddress{%
1234 Elm Street, Apt 567\\
Springfield, 98765\\
Sabah, Malaysia%
}

% -------------------------
% DECLARATION TEXT (inside the box on the Declaration page)
% -------------------------
\declarationLLM{%
I acknowledge that during the preparation of this thesis, I have utilized Large
Language Models (LLMs) as a supplementary tool to assist in various aspects of
my research and writing process. Specifically, I have used LLMs for paraphrasing
and refining my writing to improve clarity, coherence, and readability while
ensuring the original meaning remains intact. Additionally, I have leveraged LLMs
to generate general ideas and structure my thoughts more effectively, using them
as a brainstorming aid rather than as a replacement for original critical thinking.
LLMs have also been helpful in summarizing complex topics, identifying key points
from lengthy articles, and suggesting alternative ways to present information.
Furthermore, I have used LLMs to check grammar and syntax, ensuring that my
writing adheres to academic standards. In some instances, LLMs provided
guidance on citation formats and referencing best practices, although all sources
have been manually verified to maintain academic integrity. Importantly, I have
critically evaluated all AI-generated content and ensured that it aligns with my
own understanding and research findings.%
}

% -------------------------
% ACKNOWLEDGEMENT (optional: keep short, 1 page max per guide)
% If you don't want it printed, you can leave it empty.
% -------------------------
\penghargaan{First and foremost, I would like to express my sincere gratitude to my supervisor, Dr. Rodney Petrus Balandong, for his invaluable guidance, continuous support, and constructive feedback throughout the course of this research. His expertise,  patience, and dedication to academic excellence greatly contributed to the successful completion of this thesis.
I would also like to acknowledge the Faculty of Science and Technology, Universiti Malaysia Sabah, for providing the academic environment, facilities, and resources necessary to conduct this study. The support and assistance from the faculty members and administrative staff are sincerely appreciated.
My deepest appreciation is extended to my family for their unwavering encouragement, understanding, and moral support throughout my academic journey. Their constant motivation and belief in me have been a source of strength during challenging times.
I would also like to thank my friends and peers who provided helpful discussions, technical assistance, and encouragement throughout this research. Their support contributed positively to the completion of this work.
}

% -------------------------
% ABSTRACT (English) - 150 to 250 words, single paragraph, single spacing
% fin.sty will print it; you only store the text here.
% -------------------------
\abstract{Partial discharge (PD) is a critical indicator of insulation degradation in high-voltage electrical equipment, and accurate PD classification is essential for condition-based maintenance and early fault prevention. Conventional PD diagnostic approaches rely heavily on expert interpretation of phase-resolved discharge patterns, which are subjective, time-consuming, and vulnerable to noise and complex operating conditions. To address these limitations, this study investigates the use of traditional machine learning techniques for automated partial discharge classification using a structured, multi-domain feature-based framework. A publicly available partial discharge dataset was utilized, and features were extracted from time-domain, statistical, frequency-domain, and time--frequency representations to capture the physical and spectral characteristics of PD signals. Two independent feature selection strategies---Featurewiz and MLJAR AutoML---were employed to identify the most informative features while reducing dimensionality. Several traditional machine learning classifiers, including Support Vector Machine, Random Forest, Decision Tree, k-Nearest Neighbors, and Artificial Neural Network models, were trained and evaluated using stratified k-fold cross-validation. Model performance was assessed using accuracy, precision, recall, and F1-score metrics. The results show that Random Forest and Artificial Neural Network classifiers consistently achieved the strongest performance across both feature selection strategies, with classification accuracies exceeding 90\% and stable F1-scores. Despite differences in feature selection mechanisms, both approaches identified a common core of discriminative features, particularly those related to phase-resolved discharge behavior and energy distribution across frequency bands. Overall, this study confirms the practical viability of traditional machine learning approaches for robust, interpretable, and computationally efficient partial discharge classification in real-world power system monitoring applications.
}

% -------------------------
% ABSTRACT (Bahasa Malaysia) - required after English abstract
% -------------------------
\abstractMalay{Pelepasan separa (Partial Discharge, PD) merupakan petunjuk awal yang kritikal terhadap degradasi penebatan dalam peralatan elektrik voltan tinggi, dan pengelasan PD yang tepat adalah penting bagi penyelenggaraan berasaskan keadaan serta pencegahan kegagalan awal. Kaedah diagnostik PD konvensional bergantung secara meluas kepada tafsiran pakar terhadap corak pelepasan terlerai fasa, yang bersifat subjektif, memakan masa, dan terdedah kepada gangguan hingar serta keadaan operasi yang kompleks. Bagi mengatasi kekangan ini, kajian ini menilai penggunaan teknik pembelajaran mesin tradisional untuk pengelasan pelepasan separa secara automatik menggunakan rangka kerja ciri berbilang domain yang tersusun. Set data pelepasan separa yang tersedia secara umum telah digunakan, dan ciri-ciri diekstrak daripada domain masa, statistik, domain frekuensi, dan masa--frekuensi bagi menangkap ciri fizikal dan spektrum isyarat PD. Dua strategi pemilihan ciri yang bebas, iaitu Featurewiz dan MLJAR AutoML, telah diaplikasikan bagi mengenal pasti ciri paling diskriminatif serta mengurangkan dimensi ciri. Beberapa pengelas pembelajaran mesin tradisional, termasuk Support Vector Machine, Random Forest, Decision Tree, k-Nearest Neighbors, dan Artificial Neural Network, telah dilatih dan dinilai menggunakan pengesahan silang k-lipatan berstrata. Prestasi model dinilai berdasarkan metrik ketepatan, kejituan, kebolehan panggilan semula, dan skor F1. Keputusan menunjukkan bahawa pengelas Random Forest dan Artificial Neural Network secara konsisten mencapai prestasi terbaik merentas kedua-dua strategi pemilihan ciri, dengan ketepatan pengelasan melebihi 90\% dan skor F1 yang stabil. Walaupun terdapat perbezaan dalam mekanisme pemilihan ciri, kedua-dua pendekatan mengenal pasti satu set teras ciri yang signifikan, khususnya yang berkaitan dengan tingkah laku pelepasan terlerai fasa dan taburan tenaga merentasi jalur frekuensi. Secara keseluruhan, kajian ini mengesahkan kebolehlaksanaan praktikal kaedah pembelajaran mesin tradisional untuk pengelasan pelepasan separa yang teguh, boleh ditafsir, dan cekap dari segi pengiraan dalam aplikasi pemantauan sistem kuasa sebenar.
}

%   \beforepreface
%   \afterpreface

% -------------------------
% REQUIRED METADATA
% -------------------------
\title{PARTIAL DISCHARGE CLASSIFICATION USING MACHINE LEARNING}
\author{MALCOLM SABINUS}

\degree{BACHELOR OF SCIENCE WITH HONORS INDUSTRIAL PHYSICS}
\dept{INDUSTRIAL PHYSICS}
\mainsupervisor{Dr. Rodney Petrus Balandong}
\matricno{BS22110396}

% Submission date format (MONTH YEAR)
\submitdate{%
\ifcase\the\month\or
JANUARY\or FEBRUARY\or MARCH\or APRIL\or MAY\or JUNE\or
JULY\or AUGUST\or SEPTEMBER\or OCTOBER\or NOVEMBER\or DECEMBER\fi
\space \number\the\year}

\copyrightyear{\number\the\year}

% NOTE: Avoid \hline here. It's table-only and can break formatting.
% If you need separators, use plain \\ only.
\authoraddress{%
1234 Elm Street, Apt 567\\
Springfield, 98765\\
Sabah, Malaysia%
}

% -------------------------
% DECLARATION TEXT (inside the box on the Declaration page)
% -------------------------
\declarationLLM{%
I acknowledge that during the preparation of this thesis, I have utilized Large
Language Models (LLMs) as a supplementary tool to assist in various aspects of
my research and writing process. Specifically, I have used LLMs for paraphrasing
and refining my writing to improve clarity, coherence, and readability while
ensuring the original meaning remains intact. Additionally, I have leveraged LLMs
to generate general ideas and structure my thoughts more effectively, using them
as a brainstorming aid rather than as a replacement for original critical thinking.
LLMs have also been helpful in summarizing complex topics, identifying key points
from lengthy articles, and suggesting alternative ways to present information.
Furthermore, I have used LLMs to check grammar and syntax, ensuring that my
writing adheres to academic standards. In some instances, LLMs provided
guidance on citation formats and referencing best practices, although all sources
have been manually verified to maintain academic integrity. Importantly, I have
critically evaluated all AI-generated content and ensured that it aligns with my
own understanding and research findings.%
}

% -------------------------
% ACKNOWLEDGEMENT (optional: keep short, 1 page max per guide)
% If you don't want it printed, you can leave it empty.
% -------------------------
\penghargaan{First and foremost, I would like to express my sincere gratitude to my supervisor, Dr. Rodney Petrus Balandong, for his invaluable guidance, continuous support, and constructive feedback throughout the course of this research. His expertise,  patience, and dedication to academic excellence greatly contributed to the successful completion of this thesis.
I would also like to acknowledge the Faculty of Science and Technology, Universiti Malaysia Sabah, for providing the academic environment, facilities, and resources necessary to conduct this study. The support and assistance from the faculty members and administrative staff are sincerely appreciated.
My deepest appreciation is extended to my family for their unwavering encouragement, understanding, and moral support throughout my academic journey. Their constant motivation and belief in me have been a source of strength during challenging times.
I would also like to thank my friends and peers who provided helpful discussions, technical assistance, and encouragement throughout this research. Their support contributed positively to the completion of this work.
}

% -------------------------
% ABSTRACT (English) - 150 to 250 words, single paragraph, single spacing
% fin.sty will print it; you only store the text here.
% -------------------------
\abstract{Partial discharge (PD) is a critical indicator of insulation degradation in high-voltage electrical equipment, and accurate PD classification is essential for condition-based maintenance and early fault prevention. Conventional PD diagnostic approaches rely heavily on expert interpretation of phase-resolved discharge patterns, which are subjective, time-consuming, and vulnerable to noise and complex operating conditions. To address these limitations, this study investigates the use of traditional machine learning techniques for automated partial discharge classification using a structured, multi-domain feature-based framework. A publicly available partial discharge dataset was utilized, and features were extracted from time-domain, statistical, frequency-domain, and time--frequency representations to capture the physical and spectral characteristics of PD signals. Two independent feature selection strategies---Featurewiz and MLJAR AutoML---were employed to identify the most informative features while reducing dimensionality. Several traditional machine learning classifiers, including Support Vector Machine, Random Forest, Decision Tree, k-Nearest Neighbors, and Artificial Neural Network models, were trained and evaluated using stratified k-fold cross-validation. Model performance was assessed using accuracy, precision, recall, and F1-score metrics. The results show that Random Forest and Artificial Neural Network classifiers consistently achieved the strongest performance across both feature selection strategies, with classification accuracies exceeding 90\% and stable F1-scores. Despite differences in feature selection mechanisms, both approaches identified a common core of discriminative features, particularly those related to phase-resolved discharge behavior and energy distribution across frequency bands. Overall, this study confirms the practical viability of traditional machine learning approaches for robust, interpretable, and computationally efficient partial discharge classification in real-world power system monitoring applications.
}

% -------------------------
% ABSTRACT (Bahasa Malaysia) - required after English abstract
% -------------------------
\abstractMalay{Pelepasan separa (Partial Discharge, PD) merupakan petunjuk awal yang kritikal terhadap degradasi penebatan dalam peralatan elektrik voltan tinggi, dan pengelasan PD yang tepat adalah penting bagi penyelenggaraan berasaskan keadaan serta pencegahan kegagalan awal. Kaedah diagnostik PD konvensional bergantung secara meluas kepada tafsiran pakar terhadap corak pelepasan terlerai fasa, yang bersifat subjektif, memakan masa, dan terdedah kepada gangguan hingar serta keadaan operasi yang kompleks. Bagi mengatasi kekangan ini, kajian ini menilai penggunaan teknik pembelajaran mesin tradisional untuk pengelasan pelepasan separa secara automatik menggunakan rangka kerja ciri berbilang domain yang tersusun. Set data pelepasan separa yang tersedia secara umum telah digunakan, dan ciri-ciri diekstrak daripada domain masa, statistik, domain frekuensi, dan masa--frekuensi bagi menangkap ciri fizikal dan spektrum isyarat PD. Dua strategi pemilihan ciri yang bebas, iaitu Featurewiz dan MLJAR AutoML, telah diaplikasikan bagi mengenal pasti ciri paling diskriminatif serta mengurangkan dimensi ciri. Beberapa pengelas pembelajaran mesin tradisional, termasuk Support Vector Machine, Random Forest, Decision Tree, k-Nearest Neighbors, dan Artificial Neural Network, telah dilatih dan dinilai menggunakan pengesahan silang k-lipatan berstrata. Prestasi model dinilai berdasarkan metrik ketepatan, kejituan, kebolehan panggilan semula, dan skor F1. Keputusan menunjukkan bahawa pengelas Random Forest dan Artificial Neural Network secara konsisten mencapai prestasi terbaik merentas kedua-dua strategi pemilihan ciri, dengan ketepatan pengelasan melebihi 90\% dan skor F1 yang stabil. Walaupun terdapat perbezaan dalam mekanisme pemilihan ciri, kedua-dua pendekatan mengenal pasti satu set teras ciri yang signifikan, khususnya yang berkaitan dengan tingkah laku pelepasan terlerai fasa dan taburan tenaga merentasi jalur frekuensi. Secara keseluruhan, kajian ini mengesahkan kebolehlaksanaan praktikal kaedah pembelajaran mesin tradisional untuk pengelasan pelepasan separa yang teguh, boleh ditafsir, dan cekap dari segi pengiraan dalam aplikasi pemantauan sistem kuasa sebenar.
}
